\section{Методология экспериментального исследования}

\subsection{Общая схема}
Исследование проводилось в несколько этапов, соответствующих эволюции архитектурного подхода (подробно — в разделе 2.4). Финальная методология базируется на гибридной схеме «нейросетевая сегментация + детерминированная CAD-реконструкция». Основные этапы методологии:

\begin{enumerate}
    \item \textbf{Подготовка и очистка датасета}: из сырых данных (папки с STL, STEP, JSON) формируется размеченный набор для обучения классификаторов и сегментаторов.
    \item \textbf{Обучение нейросетевых моделей}: для каждой операции независимо обучаются классификатор (бинарная классификация наличия операции) и сегментатор (попиксельная/пореберная маска) на основе архитектуры MeshCNN.
    \item \textbf{Инференс и постобработка}: входная STEP-модель триангулируется, нормализуется, подаётся на обученные сети; результаты сегментации фильтруются и преобразуются в STL выделенной геометрии.
    \item \textbf{Интеграция с КОМПАС-3D}: полученный STL загружается в CAD-среду, где по нему детерминированно вычисляются параметры операции и выполняется упрощение модели.
    \item \textbf{Оценка качества}: вычисляются метрики accuracy для классификации и доля верно классифицированных рёбер для сегментации, а также визуально анализируются типовые ошибки.
\end{enumerate}

\subsection{Инструментарий и программное обеспечение}
\begin{itemize}
    \item \textbf{Язык программирования}: Python 3.9.
    \item \textbf{Библиотеки машинного обучения}: PyTorch 1.12, MeshCNN (форк от авторов), TensorBoard.
    \item \textbf{Обработка геометрии}: trimesh, numpy-stl, opencascade (через PythonOCC) для работы с STEP.
    \item \textbf{Среда выполнения инференса}: Windows 10/11, КОМПАС-3D версии 20 и выше.
    \item \textbf{Аппаратное обеспечение для обучения}: NVIDIA RTX 4090 (24 ГБ VRAM), 64 ГБ RAM, SSD 1 ТБ.
\end{itemize}

\subsection{Метрики оценки}
Для классификации использовалась стандартная \textbf{accuracy} (доля правильно предсказанных наличия операции на тестовой выборке). Для сегментации введены две метрики:
\begin{itemize}
    \item \textbf{Edge accuracy} — доля рёбер, которым присвоена правильная метка (0 — фон, 1 — операция).
    \item \textbf{Model accuracy @90} — доля моделей, в которых не менее 90\% рёбер классифицированы верно.
\end{itemize}
Такие метрики позволяют оценить как локальную точность сети, так и её способность правильно обрабатывать целые модели.

